% concatenated from main.tex

%-----------------------------------------------------------------------
%    Beginning of article.tex
%-----------------------------------------------------------------------
%
%    This is an AMS-LaTeX sample proceedings article file for use with
%    the amsproc document class and author packages based on amsproc.
%
%    Replace amsproc by the document class name for the target series,
%    e.g. pspum-l.
%
\documentclass{amsproc}
\usepackage{xcolor}
\usepackage{amssymb}
\newtheorem{theorem}{Theorem}[section]
\newtheorem{lemma}[theorem]{Lemma}
\newtheorem{op}[theorem]{Open}
\newtheorem{proposition}[theorem]{Proposition}
\newtheorem{corollary}[theorem]{Corollary}

\theoremstyle{definition}
\newtheorem{definition}[theorem]{Definition}
\newtheorem{example}[theorem]{Example}
\newtheorem{xca}[theorem]{Exercise}

\theoremstyle{remark}
\newtheorem{remark}[theorem]{Remark}
\newtheorem{note}[theorem]{Note}
\numberwithin{equation}{section}
\usepackage{epsfig} 
\usepackage{epstopdf} 
%Lattice operations
\newcommand{\jj}{\vee}% join
\newcommand{\mm}{\wedge}% meet
\newcommand{\JJ}{\bigvee}% big join
\newcommand{\MM}{\bigwedge}% big meet
\newcommand{\JJm}[2]{\JJ(\,#1\mid#2\,)}% big join with a middle
\newcommand{\MMm}[2]{\MM(\,#1\mid#2\,)}% big meet with a middle


%Set operation
\newcommand{\uu}{\cup}% union
\newcommand{\ii}{\cap}% intersection
\newcommand{\UU}{\bigcup}% big union
\newcommand{\II}{\bigcap}% big intersection
\newcommand{\UUm}[2]{\UU\{\,#1\mid#2\,\}}% big union with a middle
\newcommand{\IIm}[2]{\II\{\,#1\mid#2\,\}}% big intersection with a middle

%Sets
\newcommand{\ci}{\subseteq}% contained in with equality
\newcommand{\nc}{\nsubseteq}% not \ci
\newcommand{\sci}{\subset}% strictly contained in
\newcommand{\nci}{\nc}% not \ci
\newcommand{\ce}{\supseteq}% containing with equality
\newcommand{\sce}{\supset}
\newcommand{\nce}{\nsupseteq}% not \ce
\newcommand{\nin}{\notin}% not \in
\newcommand{\es}{\emptyset}% the empty set
\newcommand{\set}[1]{\{#1\}}% set
\newcommand{\setm}[2]{\{\,#1\mid#2\,\}}% set with a middle

%Partial ordering
\newcommand{\nle}{\nleq}% not \leq

%Greek letters
\newcommand{\ga}{\alpha}
\newcommand{\gb}{\beta}
\newcommand{\gc}{\chi}
\newcommand{\gd}{\delta}
\renewcommand{\gg}{\gamma}% old use >>
\newcommand{\gh}{\eta}
\newcommand{\gi}{\iota}
\newcommand{\gk}{\kappa}
\newcommand{\gl}{\lambda}
\newcommand{\go}{\omega}
\newcommand{\gp}{\pi}
\newcommand{\gq}{\theta}
\newcommand{\gs}{\sigma}
\newcommand{\gt}{\tau}
\newcommand{\gx}{\xi}
\newcommand{\gy}{\psi}
\newcommand{\gz}{\zeta}
\newcommand{\vp}{\varphi}

\newcommand{\gG}{\Gamma}
\newcommand{\gD}{\Delta}
\newcommand{\gF}{\Phi}
\newcommand{\gL}{\Lambda}
\newcommand{\gO}{\Omega}
\newcommand{\gP}{\Pi}
\newcommand{\gQ}{\mathfrak{\nu}}
\newcommand{\gS}{\Sigma}
\newcommand{\gX}{\Xi}
\newcommand{\gY}{\Psi}

%Font command
\newcommand{\tbf}{\textbf}% text bold
\newcommand{\tit}{\textit}% text italic

\newcommand{\mbf}{\mathbf}% math bold
\newcommand{\B}{\boldsymbol}% Bold math symbol, use as \B{a}
\newcommand{\C}[1]{\mathcal{#1}}% Euler Script - only caps, use as \C{A}
\newcommand{\D}[1]{\mathbb{#1}}% Doubled - blackboard bold - only caps, uas as \D{A}
\newcommand{\F}[1]{\mathfrak{#1}}% Fraktur, use as \F{a}

%Miscellaneous
\newcommand{\te}{\text}% same as \mathrm command.
\newcommand{\tei}{\textit}
\newcommand{\im}{\implies}
\newcommand{\ti}{\times}
\newcommand{\exi}{\ \exists \ }
\newcommand{\ther}{\therefore}
\newcommand{\be}{\because}
\newcommand{\ep}{\epsilon}
\newcommand{\la}{\langle}
\newcommand{\ra}{\rangle}
\newcommand{\ol}{\overline}
\newcommand{\ul}{\underline}
\newcommand{\q}{\quad}% spacing
\newcommand{\qq}{\qquad}% morepspacing
\newcommand{\fa}{\ \forall \, }
\newcommand{\nd}{\noindent}
\newcommand{\bs}{\backslash}
\newcommand{\pa}{\partial}
\newcommand{\para}{\parallel}
\newcommand{\sm}{\setminus}
\newcommand{\tl}{\tilde}


%    Absolute value notation
\newcommand{\abs}[1]{\lvert#1\rvert}

%    Blank box placeholder for figures (to avoid requiring any
%    particular graphics capabilities for printing this document).
\newcommand{\blankbox}[2]{%
  \parbox{\columnwidth}{\centering
%    Set fboxsep to 0 so that the actual size of the box will match the
%    given measurements more closely.
    \setlength{\fboxsep}{0pt}%
    \fbox{\raisebox{0pt}[#2]{\hspace{#1}}}%
  }%
}

\begin{document}

\title[Weaker quantization dimension results for self-similar measures]{Weaker quantization dimension results for self-similar measures}

%    Information for first author

%    \thanks will become a 1st page footnote.


%    Information for the second author
\author{Saurabh Verma}
\address{Department of Applied Sciences, IIIT Allahabad, Prayagraj, India 211015}
\email{saurabhverma@iiita.ac.in}
\thanks{The first author is supported by the SEED Grant Project of IIIT Allahabad, India. The second author is financially supported by the Ministry of Education, India.}

\author{Shivam Dubey}
%    Address of record for the research reported here
\address{Department of Applied Sciences, IIIT Allahabad, Prayagraj, India 211015}
%    Current address
%
\email{rss2022509@iiita.ac.in}





\subjclass[2020]{Primary 28A80; Secondary 28A78}
\date{January 1, 1994 and, in revised form, June 22, 1994.}

% \dedicatory{This paper is dedicated to our advisors.}

\keywords{Separation condition, self-similar sets, invariant measures, Lower Assouad dimension, Quantization dimension}

\begin{abstract}
In this paper, we investigate the quantization dimension of self-similar measures, particularly when the IFS does not satisfy the separation condition, but the sub-IFS at some level satisfies the separation condition. Further, we study the approximation of the space of Borel probability measures $\Omega(\mathbb{R}^m)$ with respect to the geometric mean error, i.e., the quantization dimension of order zero.

\end{abstract}

\maketitle

\section{Introduction}
An important factor in the study of fractal geometry is the fractal dimension. It measures a fractal's scaling across various magnifications to provide information on geometric complexity; it translates \textit{irregularity} into an exact numerical quantity, the \textit{fractal dimension}, that interpolates between traditional geometric dimensions. Since the irregularity exhibited by the fractal is not always homogeneous, a single notion of the fractal dimension is not sufficient. There are various dimensions that have been introduced in the literature to capture the different aspects of the complexity, including the Hausdorff dimension $\dim_\mathcal{H}(\cdot),$ Minkowski–Bouligand dimension (box-counting dimension) $\dim_B(\cdot),$ packing dimension $\dim_P(\cdot),$ Assouad dimension $\dim_A(\cdot),$ lower Assouad dimension (lower dimension) $\dim_L(\cdot),$ quantization dimension $\dim_r(\cdot), ~ L^q-$ dimension, information dimension, correlation dimension, entropy dimension, etc. (see \cite{Fal, Fraser, KNZ1, Mat3, Sh1, Zador, KPot}). Some dimensions, such as the Hausdorff and box dimensions, capture global scaling complexity, while others, such as the Assouad and lower dimensions, describe local extremal behaviour. In the case of fractal measures, global notions such as the information dimension and the $L^q$-dimension, as well as local dimensions such as the local dimension and the multifractal spectrum, provide insight at the local level. 
The concept of \textit{quantization} was first derived in information theory and signal processing, where it refers to the approximation of a continuous signal by a finite set of discrete code points. Zadore \cite{Zador} traces this idea and derived the asymptotic formulas for the optimal quantization error at high rates of a probability distribution. In this context, a natural question is to determine the rate of convergence of the optimal-$n$ approximation error as $n$ approaches infinity. Building on these ideas, Graf and Luschgy \cite{GL1} developed a measure-theoretic framework for quantization of probability distribution and formalized the notion of \textit{quantization coefficient} and \textit{quantization dimension}. In general, calculating the fractal dimension exactly is a tedious task. It becomes easier when the cylinders have no or minimal overlap. There are various separation properties in the literature, such as \textit{strong separation condition (SSC)}, \textit{open set condition (OSC)}, \textit{strong open set condition (SOSC)}, and \textit{exponential separation condition (ESC)}.


\section{Preliminaries} \label{Sec2}
In this section, we recall some definitions and basic results for our paper. For any $F\subseteq\mathbb{R}^m,$ we denote $|F|=\sup_{x,y\in F}\|x-y\|_2$ and $\text{Card}(F)$ as diameter and cardinality of $F,$ respectively. \\
Let $\mathfrak{K}(\mathbb{R}^m)$ denote the collection of all compact non-empty subsets of $\mathbb{R}^m$ with the Euclidean norm $\|\cdot\|_2$. The Hausdorff distance $\mathfrak{H}$ between any $E,F\in\mathfrak{K}(\mathbb{R}^m)$ is elaborated as 
% \begin{equation*}
%     \mathfrak{H}(A,B):=\max\big\{\max_{a\in A}D(a,B),\max_{b\in B}D(b,A)\big\},
% \end{equation*}
% where $D(a,B)=\min_{b\in B}|a-b|.$ It is also defined as
$$\mathfrak{H}_d(E,F):=\inf\{\epsilon>0: E\subset F_\epsilon \text{ and } F\subset E_\epsilon\},$$ where $E_\epsilon$ denotes the $\epsilon$-neighbourhood of $E.$
% \begin{definition}\cite{Fraser}
%     Let $E$ be a non-empty subset of $\mathbb{R}^m.$ The lower dimension of $E$ is defined as
%     \begin{align*}
%        \dim_{L}(E):=\sup&\Big\{\beta: \exists ~M>0 \text{ such that,} \forall ~0<r<R\le|E| \text{ and }x\in E, \\& \hspace{4cm}N_r(\mathcal{B}_R(x)\cap E)\ge M\Big(\frac{R}{r}\Big)^\beta\Big\},
%     \end{align*}
% where $\mathcal{B}_R(x)=\{y\in\mathbb{R}^m:\|y-x\|_2\le R\}$ and $N_r(\mathcal{B}_R(x)\cap E)$ is the smallest number of open set required for a $r$-cover of $\mathcal{B}_R(x)\cap E.$    
% \end{definition}
% The properties of lower dimensions are discussed in the book \cite{Fraser}. It is noted that the lower dimension satisfies stability only under closure and bi-Lipschitz mappings.
% \begin{definition}\cite{Fraser}
%  Let $\mu$ be a Borel probability measure on $\mathbb{R}^m$ and $\text{supp}(\mu)$ denote the support of $\mu.$ Then 
 % the Assouad dimension of measure $\mathfrak$ is defined by 
 % \begin{align*}
 %      \dim_A(\mathfrak)=\sup\Bigg\{&\beta\ge0:\text{ there exists a constant } M>0 \text{ such that},\text{ for all }\\ &0<r<R<|\text{supp}(\mathfrak)|\text{ and }x\in\text{supp}(\mathfrak),~\frac{\mathfrak(\mathcal{B}_R(x))}{\mathfrak(\mathcal{B}_R(x))}\le M\Big(\frac{R}{r}\Big)^\beta\Bigg\},
 % \end{align*}
 % and provided $|\text{supp}| >0,$ and 
%  the lower dimension of $\mu$ is defined by
%  \begin{align*}
% \dim_L(\mu)=\sup\Bigg\{&\beta\ge0:\exists ~M>0 \text{ such that,} \forall ~0<r<R\le|E| \text{ and }x\in \text{supp}(\mu),\\& \hspace{5.8cm}\frac{\mu(\mathcal{B}_R(x))}{\mu(\mathcal{B}_r(x))}\ge M\Big(\frac{R}{r}\Big)^\beta\Bigg\},
%  \end{align*}
% otherwise it is $0.$
% \end{definition}

% We would like to mention here that if a measure $\mu$ is doubling and fully supported on a closed set $F,$ then $\dim_L(\mu)\le \dim_L(F),$
% for details, see \cite[Lemma 4.1.2]{Fraser}.
% \begin{definition}\cite{Fal}
% Let $E$ be a subset of $\mathbb{R}^m$ and $r$ be any non-negative number. Then the $r$-dimensional Hausdorff measure of $E$ is defined as $$\mathcal{H}^r(E)=\lim_{\delta\to 0^+}\mathcal{H}_{\delta}^r(E),$$
% where $\mathcal{H}_{\delta}^r(E)=\inf\{\sum_{i=1}^{\infty}(diam(U_i))^r:~ \{U_i\} \text{ is a } \delta\text{-cover of } E \}$ and $diam(U_i)=\sup_{x,y\in U_i}\|x-y\|_2.$ The Hausdorff dimension of $E$ is defined as $$\dim_{\mathcal{H}}E=\inf\{r\geq0:~\mathcal{H}^r(E)=0\}.$$
% \end{definition}
% \begin{definition}\cite{Fraser}
% Let $E$ be a non-empty subset of $\mathbb{R}^m.$ The Assouad dimension of $E$ is defined as
% \begin{align*}
%    \dim_{\mathcal{A}}E=\inf\Big\{\beta: \exists \text{ a constant }&C>0\text{ such that }\forall~ 0<r<R\text{ and }x\in E,\\&N_r(\mathcal{B}_R(x)\cap E)\le C\Big(\frac{R}{r}\Big)^\beta\Big\}, 
% \end{align*}
% where $\mathcal{B}_R(x)=\{y\in\mathbb{R}^m:\|y-x\|_2\le R\}$ and $N_r(\mathcal{B}_R(x)\cap E)$ is the smallest number of open set required for a $r$-cover of $\mathcal{B}_R(x)\cap E.$    
% \end{definition}
% \begin{example}\label{eg1}
% The Assouad dimension of any finite subset of $\mathbb{R}^m$ is zero.    
% \end{example}
% We now mention properties of the Assouad dimension in the following theorem, used throughout the article.
% \begin{theorem}\cite[Lemma 2.4.1, Lemma 2.4.2]{Fraser}\label{th1}
% \begin{itemize}
%     \item[-] The Assouad dimension is monotone, finitely stable, stable under closure, and satisfies the open set property. 
%     \item[-] Let $E$ be a non-empty subset of $\mathbb{R}^m,$ and $T:\mathbb{R}^m\rightarrow\mathbb{R}^m$ be a bi-Lipschitz map. Then $\dim_\mathcal{A}T(E)=\dim_\mathcal{A}E.$
% \end{itemize}   
% \end{theorem}
% Let $\Omega(\mathbb{R}^m)$ be the set of probabilty measure on $\mathbb{R}^m$ with the metric defined as
% $$d_L(\mathfrak,\mathfrak{\nu})=\sup_{Lip(f)\leq1}\Big\{\Big|\int_{\mathbb{R}^m}fd\mathfrak-\int_{\mathbb{R}^m}fd\mathfrak{\nu}\Big|\Big\},$$
% where $f$ is a Lipschitz function with the Lipschitz constant equal to $1$. 
% Let $\mathfrak,\mathfrak{\nu}\in\Omega(\mathbb{R}^m).$ Then the convolution $\mathfrak*\mathfrak{\nu}$ is given as 
%     $$\mathfrak*\mathfrak{\nu}(A)=\int_{\mathbb{R}^m}\mathfrak(A-x) d\mathfrak{\nu}(x),\hspace{.4cm} \mbox{for all Borel set } A \in \mathbb{R}^m.$$
%       Equivalently $\mathfrak*\mathfrak{\nu}$ is the unique Borel probability measure satisfying 
%       \begin{equation*}
% \int_{\mathbb{R}^m}f(x)\mathfrak*\mathfrak{\nu}(A)=\int_{\mathbb{R}^m}\int_{\mathbb{R}^m}f(x+y)d\mathfrak(x)d\mathfrak{\nu}(y),\vspace{.4cm} \forall f\in C_b(\mathbb{R}^m)\end{equation*}
%  where, $C_b(\mathbb{R}^m)$ denotes the set of all continuous and bounded functions on $\mathbb{R}^m$. \\
%  If $\mathfrak=\sum_{j=1}^nc_j\delta_{x_j}$ by above equation,
% \begin{eqnarray}\label{convolusion}
% \int_{\mathbb{R}^m}f(x)\mathfrak*\mathfrak{\nu}(A)&=&\int_{\mathbb{R}^m}\int_{\mathbb{R}^m}f(x+y)d(\sum_{j=1}^nc_j\delta_{x_j})d\mathfrak{\nu}(y)\nonumber  \\
% &=& \int_{\mathbb{R}^m}\sum_{j=1}^nc_jf(x_i+y)d\mathfrak{\nu}(y)
% \end{eqnarray}\\
% We know that (i) $\mathfrak*\mathfrak{\nu}(A)=\mathfrak{\nu}*\mathfrak(A)$ and (ii)  for any $x_0\in\mathbb{R}^m,$ we have  $\mathfrak*\delta_{x_0}(A)=\mathfrak(A-x_0).$
% $$\mathfrak*\mathfrak{\nu}(A)=\int_{\mathbb{R}^m}\int_{\mathbb{R}^m}1_A(x+y)d\mathfrak(x)d\mathfrak{\nu}(y),$$ where $A$ is a measurable subset of $\mathbb{R}^m.$
% \item Let $x_0\in\mathbb{R}^m.$ Then $\mathfrak*\delta_{x_0}(A)=\mathfrak(A-x_0),$ where 
% %%%%%%%
\begin{definition}\cite{GL1}
    Let $\mu$ be a Borel probability measure on $\mathbb{R}^m,~r \in (0, +\infty)$ and $n \in \mathbb{N}.$  Then, the $n$th quantization error of order $r$ of $\mu$ is stated as:
    \[V_{n, r}(\mu):=\text{inf} \Big\{\int d(x, A)^r d\mu(x): A \subset \mathbb{R}^m, \ \text{Card}(A) \leq n\Big\},\]
    where $d(x,A)$ denotes the distance of point $x \in \mathbb{R}^m$ to the set $A$ with respect to the Euclidean norm $\|\cdot\|_2$ on $\mathbb{R}^m.$
    The quantization dimension of order $r$ of $\mu$ is defined by:
    $$ D_r(\mu):= \lim_{n\to \infty} \frac{r \log n}{- \log(V_{n,r}(\mu))}.$$
    If the limit exists, otherwise, we define the lower and upper quantisation dimensions of order $r$ of $\mu$ by taking the limit inferior and limit superior of the sequence, respectively.
\end{definition}
\begin{definition}\cite{GL1}
    For a given $s>0,$ the $s$-dimensional lower and upper quantization coefficient of order $r$ for $\mu$ is defined by:
    $$\liminf_{n \to \infty} n^{\frac{r}{s}}V_{n,r}(\mu),~ \text{and}~ \limsup_{n \to \infty}  n^{\frac{r}{s}}V_{n,r}(\mu),$$
    respectively.
\end{definition}
Let $\Omega(\mathbb{R}^m)$ denote the set of all Borel probability measures on the Euclidean space $(\mathbb{R}^m,\|.\|_2)$. Then,
\begin{equation*}
\begin{aligned}
d_L (\mathfrak{\mu}, \mathfrak{\nu}) :=\sup_{h\in\te{Lip}_1(\mathbb{R}^m)} \Big\{\Big|\int_{\mathbb{R}^m} h d\mu -\int_{\mathbb{R}^m} hd\mathfrak{\nu}\Big|\Big\}, \ \ \quad \mu, \mathfrak{\nu} \in \Omega(\mathbb{R}^m) , \end{aligned}
\end{equation*}
defines a metric on $\Omega(\mathbb{R}^m)$, where $\te{Lip}_1(\mathbb{R}^m)=\{h:\mathbb{R}^m\rightarrow \mathbb{R}: h\te{ is Lipschitz function}\\ \te{with Lipschitz constant }\le1\}$ (see \cite{Mat,Parth}). If $(X,d)$ is a compact metric space, then
$(\Omega(X), d_L)$ is a compact metric space (see \cite[Theorem~5.1]{B}), also note that if $X$ is separable and complete then the space $\omega(\mathbb{R}^m)$ is complete.
% Again, we know in the weak topology on $\Omega(X)$,
% \[\mathfrak_n \to \mathfrak  \Longleftrightarrow \int_X f d\mathfrak_n - \int_X f d\mathfrak \to 0 \te{ for all } f \in \C C(X),\]
% where $\C C (X) : =\set{ f : X \to \D R : f \te{ is continuous}}$. It is known that the $d_L$-topology and the weak topology coincide on the space of probabilities with compact support (see \cite{Mat}). 
%In our case, all measures are compactly supported.
\par

%We refer to \cite{Shaw} for the definitions of generic (something $G_{\delta}$-dense set) and prevalence.


% \begin{definition}
% The Fourier transform of $\mathfrak \in \Omega(\mathbb{R}^m)$ is defined as 
% \[
% \hat{\mathfrak}(\xi)= \int e^{-2\pi i \xi .x} d\mathfrak(x),~~~~ \xi \in \mathbb{R}^m.
% \]
% \end{definition}



\begin{definition}
Let $\mu, \mathfrak{\nu} \in \Omega(\mathbb{R}^m).$ The convolution of $\mu$ and $\mathfrak{\nu}$ is denoted by $\mu *\mathfrak{\nu}$ and defined by the push-forward of the product measure $\mu \times \mathfrak{\nu}$ under the mapping $(x,y)\to x+y.$ That is, for any $E \subset \mathbb{R}^m$, we have $$\mu * \mathfrak{\nu} (E)= (\mu \times \mathfrak{\nu} )\big(\{(x,y) \in \mathbb{R}^m\times \mathbb{R}^m: x+y \in E\}\big).$$
\end{definition}
\begin{definition}
Let $\mu \in \Omega(\mathbb{R}^m)$ and $x \in \mathbb{R}^m.$ The translation of $\mu$ by $x$ is defined as 
\[
(\mu+x)(E)=\mu(E+x),~~~E \subset \mathbb{R}^m.
\]
\end{definition}

Let us define
\[
\mathcal{L}(\mathbb{R}^m)=\Big\{\mu 
\in \Omega(\mathbb{R}^m): \mu= \sum_{i=1}^k a_i \delta_{x_i},~x_i \in \mathbb{R}^m, ~a_i > 0, \sum_{i=1}^k a_i=1,~~ k \in \mathbb{N} \Big\}.
\]
Here $\delta_x$ denotes the Dirac measure supported at $x \in \mathbb{R}^m.$ It is observed that $\mathcal{L}(\mathbb{R}^m)$ is dense in $\Omega(\mathbb{R}^m)$ w.r.t. $d_L$ on $\Omega(\mathbb{R}^m),$ see \cite{Bill}. The existence of an invariant measure supported on the attractor of the associated iterated function system is shown in \cite{H}. Feng et al. \cite{Feng} studied the convolutions of equicontractive self-similar measures on $\mathbb{R}.$ 
% For computational purposes, let us note the following:
% \begin{itemize}
% \item Let $f \in \mathcal{C}_0(\mathbb{R}^m),$ where $\mathcal{C}_0(\mathbb{R}^m)$ denotes the set of compactly supported and continuous functions on $\mathbb{R}^m.$ Then 
% \[
% \int f d(\mathfrak * \lambda)= \int \int f(x+y)d\mathfrak(x) d\lambda(y).
% \]
% \item Let $ \mathfrak \in \Omega(\mathbb{R}^m)$ and $x \in \mathbb{R}^m.$ Then $\mathfrak * \delta_x= \mathfrak-x,$ and 
% \[
% \int f d(\mathfrak *\delta_x)= \int f(y+x) d\mathfrak(y)=\int f d(\mathfrak-x), ~~\forall ~f \in \mathcal{C}_0(\mathbb{R}^m).
% \]
% \end{itemize}
% We know that $\hat{\delta}_{\boldsymbol{0}}=1,$ where $\boldsymbol{0} \in \mathbb{R}^m.$ It is interesting to see that $$D_r(\delta_{\boldsymbol{0}})=0~ \text{and}~\dim_\mathcal{H}(Gr(\hat{\delta}_{\boldsymbol{0}}))=m,$$
% where $Gr(g):=\{(x,g(x))\in X \times Y: x \in X\}$ denotes the graph of a function $g:X \to Y.$
% By routine calculations, we also have a more general result, that is,
% $$D_r( \mathfrak)=0~ \text{and}~\dim_\mathcal{H}(Gr(\hat{\mathfrak}))=m, ~~~\forall ~\mathfrak \in \mathcal{L}(\mathbb{R}^m).$$
% It is natural to ask a similar question in general.
% \begin{remark}
% Let $f \in \mathcal{C}[0,1].$ A function $g \in  \mathcal{C}[0,1]$ is said to be multiplicative inverse of $f$ if $f(x)g(x)=1$ for every $x \in [0,1].$ We have $$ \dim_\mathcal{H}(Gr(g))=\dim_\mathcal{H}(Gr(f)).$$ To prove this, we define a bi-Lipschitz mapping $T: Gr(f) \to Gr(g)$ by $$ T(x,f(x))= \Big(x, \frac{1}{f(x)}\Big).$$ 
% \end{remark}
% In view of the above remark, we believe (not proved) that for $\mathfrak \in \Omega(\mathbb{R}^m)$, the following holds
% $$ D_r(\mathfrak{\nu})=D_r(\mathfrak),~~~ ~~\text{where}~~~\mathfrak * \mathfrak{\nu} =\delta_{\boldsymbol{0}}.$$ 
We define essential supremum norm $\|.\|_{\mathcal{L}^{\infty}}$ of a measurable function $g$ with respect to a measure $\mu$ by 
\[
 \|g\|_{\mathcal{L}^{\infty}}= \inf \{C\ge 0 : |f(x)| \le C ~\text{for $\mu$-almost every $x$} \}.
\]
Let $\mu ,\lambda \in \Omega( \mathbb{R}^m)$ be such that $\mu$ is absolutely continuous with respect to $\lambda.$ Then by Radon-Nikodym theorem, there exists a measurable function $g: \mathbb{R}^m \to [0, \infty) $ such that $$ \mu(A)= \int_A g d\lambda,~~~\text{for any Borel set}~A \subset \mathbb{R}^m.$$
Further, the function $g$ is called the Radon-Nikodym derivative or density of $\mu$ with respect to $\lambda$, and is denoted by $\frac{ d\mu}{ d\lambda}.$
We call $\mu$ has $\mathcal{L}^{\infty}$-density if it is absolutely continuous with respect to $\lambda$, and density function $g$ satisfying
 $\|g\|_{\mathcal{L}^{\infty}} < \infty.$  
% \begin{definition}\cite{Sh1}
%     Let $\mathfrak$ be a Borel probability measure on $\mathbb{R}$ with bounded support, then for $q \in (1, \infty)$ the $L^q$ spectrum of $\mathfrak$ is given by 
%     $$\tau(\mathfrak,q) := \liminf_{m \to \infty} -\frac{log \sum_{I \in \mathcal{D}_m} \mathfrak(I)^q}{m},$$
%     and the $L^q$ dimension of $\mathfrak$ is given by 
% $$D(\mathfrak,q):= \frac{\tau(\mathfrak,q)}{q-1},$$
% where for each $m \in \mathbb{N},$~$\mathcal{D}_m$ denotes the family of dyadic intervals of the form $\{[j 2^{-m},(j+1) 2^{-m}:~  j \in \mathbb{Z}\}.$
% \end{definition}
% Some of the properties, such as the dimension of the convolution of measures $L^q$ and
\par
We are now ready to prove the upcoming result. Here and throughout the paper, we assume that $D_r(\mathfrak)$ exists whenever it occurs.

% \begin{definition}
%     Let $\mathfrak$ be a Borel probability measure on $\mathbb{R^d}$ then the \textit{Assouad dimension} $\dim_A \mathfrak $ of $\mathfrak$ is defined by:
%     \begin{align*}
%         \dim_A (\mathfrak) := \inf\bigg\{s \ge 0: & \text{there exists $C>0$ such that, for all} 0<r<R< |supp(\mathfrak)|\\& \text{and $x \in supp(\mathfrak)$}, \frac{\mathfrak(\mathcal{B}_R(x))}{\mathfrak(\mathcal{B}_R(x))} \le C \bigg(\frac{R}{r}\bigg)^s \bigg\}.
%     \end{align*}
%     If $|supp(\mathfrak)|>0$ then, the \textit{lower dimension} $\dim_L (\mathfrak)$ of $\mathfrak$ is defined by:
%      \begin{align*}
%         \dim_A (\mathfrak) := \inf\bigg\{s \ge 0: & \text{there exists $C>0$ such that, for all} 0<r<R< |supp(\mathfrak)|\\& \text{and $x \in supp(\mathfrak)$}, \frac{\mathfrak(\mathcal{B}_R(x))}{\mathfrak(\mathcal{B}_R(x))} \ge C \bigg(\frac{R}{r}\bigg)^s \bigg\}.
%     \end{align*}
% \end{definition}
Let $\mathcal{I}:=(\{f_i\}_{i=1}^N, \{\rho_i\}_{i=1}^N)$ be a WIFS and $\mu$ is the corresponding self-similar measure supported on self-similar set $A.$
We denote set of all infinite sequences in $N$-symbols $\Lambda:=\{1,2,\ldots,N\}$ by $\Lambda^\infty$, and  $$\Lambda^n:= \{\xi= \xi_1 \xi_2\ldots \xi_n | \xi_i = 1,2,\ldots ,N\},~~ \Lambda^*= \bigcup_{n \in \mathbb{N}} \Lambda^n.$$
Now we define a sub-IFS at some $n$th level. Let $\xi \in \Lambda^*$ then $\mathcal{I}_n:= \{f_{\eta \xi}: \eta \in {\Upsilon}^n\}$ be a sub-IFS of the IFS $\mathcal{I}$ at $n$th level, where ${\Upsilon}^n \subseteq \Lambda^n.$ Further, the sub-IFS $\mathcal{I}_n$ with the probability vector $\big\{\frac{\rho_{\eta \xi}}{\sum_{\eta \in {\Upsilon}^n}} \rho_{\eta \xi} : \eta \in {\Upsilon}^n \big\}$ is called the sub-WIFS of the WIFS $\mathcal{I}$ at $n$th level.
\begin{note}
    Throughout the paper, we will consider $\Upsilon^n \subsetneq \Lambda^n$ in the construction of sub-IFS at the $n$-th level; otherwise, the self-similar measure $\mu_n$ associated with the sub-IFS will coincide with the self-similar measure $\mu$. 
\end{note}
\section{Main Results}
% \begin{remark}[?]
%  Let $\{\Phi_i\}_{i=1}^N$ be an IFS, and for some ${\xi} \in \Lambda^*$ the corresponding sub-IFS $\{\Phi_{ \eta \xi}: \eta \in \Lambda\}$ satisfies the SSC (OSC), then the sub-IFS $\{\Phi_{ \eta \xi} : \eta \in \Upsilon^n\}$ also satisfies the SSC (OSC), for each $n \in \mathbb{N}.$
% \end{remark} 


\begin{theorem}\label{Lem24}
    Let $\mu$ be a self-similar measure associated with a WIFS $(\{f_i\}_{i=1}^N; \\ \rho_1 ,\rho_2,\ldots,\rho_N)$ and $\mu_n$ be the self-similar measure associated with sub-WIFS $\mathcal{I}_n$ at $n$-th level. Then, $\mu_n$ is absolutely continuous with respect to $\mu$ i.e., $\mu_n << \mu.$
\end{theorem}
\begin{proof}
    Recall that, for a given probability vector $\{\varsigma_{\eta}= \frac{\rho_{ \eta \xi}}{ \sum_{\eta \in \Upsilon^n} \rho_{\eta \xi}}: \eta \in \Upsilon^n\},$ the map $\mathfrak{P}: \Omega(\mathbb{R}^m) \to \Omega(\mathbb{R}^m)$ defined by $\mathfrak{P}(\lambda)= \sum_{\eta \in \Upsilon^n} \varsigma_{\eta} \lambda \circ f_{\eta}^{-1}$ is contraction mapping with respect to the Monge-Kantorovich metric and by Banach contraction theorem it has a unique fixed point $\mu_n$. Also, independent of the choice of $\lambda \in \Omega(\mathbb{R}^m)$, the sequence $\mathfrak{P}^k(\lambda)$ converges to $\mu_n$ as $k$ approaches to infinity. Now, in order to prove that $\mu << \mu,$ first we show that if $\mu(A)=0$ for Borel set $A \subset \mathbb{R}^m,$ then $\mathfrak{P}^k(\mu)(A)= 0$ for all $k \in \mathbb{N}.$\\
    Consider, 
\begin{equation}\label{eq24}
\begin{aligned}
(\mathfrak{P} \circ \mathfrak{P})(\mu) = \mathfrak{P}\bigg(\sum_{\eta \in \Upsilon^n} \varsigma_{\eta} \mu \circ f_{\eta}^{-1} \bigg) = \sum_{\eta^1, \eta^2 \in \Upsilon^n}\varsigma_{\eta^1 \eta^2} \mu \circ f_{\eta^2 \eta^1}^{-1} 
\end{aligned}
\end{equation}
    Similarly, for $k \in \mathbb{N},$ we have 
    \[  \mathfrak{P}^k(\mu)= \sum_{\eta^1,\eta^2,\ldots,\eta^k \in \Upsilon^n} \varsigma_{\eta^1\eta^2\ldots\eta^k} \mu \circ f_{\eta^k \eta^{k-1}\ldots \eta^1}^{-1}.\]
    Since $\mu$ satisfies the self-similar equation $\mu= \sum_{i \in \Lambda} \rho_i \mu \circ f_{i}^{-1},$ we have 
    \[\mu= \sum_{i,j \in \Lambda} \rho_{ij} \mu \circ f_{ji}^{-1}.\]
    In $\zeta \in \Lambda^*$ is a word of length $l \in \mathbb{N}$ in the sub-IFS, then repeating the above process $l+n+k$ times we have 
    \begin{equation}\label{eq25}
    \begin{aligned}
        \mu(A) = \sum_{i^1,i^2,\ldots,i^{l+n+k} \in \Lambda} \rho_{i^1i^2\ldots i^{l+n+k}} \mu \circ f_{i^1i^2\ldots i^{l+n+k}}^{-1}(A)=0.
         \end{aligned}
    \end{equation}
    \end{proof}  
  A direct comparison of Equations \ref{eq24} and \ref{eq25} shows that Equation \ref{eq24} can be derived from Equation \ref{eq25} by omitting certain terms in its summation. Also, utilising the fact that the composition of similarity mapping is again a similarity mapping and $\mu$ is a Borel probability measure, we are able to establish our claim.
  Next, we claim that $\mathfrak{P}^k(\mu)$ converges setwise to $\mu_n$. Setting $\lambda = \mu$ and following the notation in \cite[Theorem 1]{XYZL}, the sequence $\{\mathfrak{P}^k(\mu)\}_{k \in \mathbb{N}}$ satisfies its hypotheses. Hence, $\mathfrak{P}^k(\mu)(A) \to \mu_n(A)$ for all Borel sets $A$, completing the proof.
\begin{corollary}\label{cor24}
    Let $\mu$ be a self-similar measure associated with a WIFS $(\{f_i\}_{i=1}^N\\; \rho_1,\rho_2,\ldots,\rho_N)$ and $\mu_n$ be the sub-IFS at $n$-th level. Then, $\mu_n$ is absolutely continuous with respect to $\mu,$ i.e., $\mu_n << \mu$ with an $L_\infty$ density.
\end{corollary}
\begin{proof}
    In order to prove the absolute continuity of measure $\mu_n$ with respect to measure $\mu$ with an $L_\infty$ density, it is sufficient to show that $\mu_n(A) \le M \mu(A)$ for some fixed constant $M>0$ and for all Borel measurable sets $A$. By Lemma \ref{Lem24} the sequence $\mathfrak{P}^k(\mu)$ converges setwise to $\mu_n$ and $\mathfrak{P}^k(\mu) \le \mu$ for all $k \in \mathbb{N}$. This implies that 
    $$\mu_n(A)= \lim_{k \to \infty} \mathfrak{P}^k(\mu(A)) \le \mu(A).$$
    This completes the assertion.
\end{proof}
  

\begin{proposition}
 Let $\mu$ be a self-similar measure associated with an IFS and $\mu_n$ be the self-similar measure corresponding to a sub-IFS at the $n$-th level. Then, $D_r(\mu_n) \le D_r(\mu)$ and $\dim_L(\mu_n) \le \dim_L(\mu).$
%provided the sub-IFS at the $n$-th level satisfies OSC and SSC, respectively.
\end{proposition}
\begin{proof}
In Theorem \ref{Lem24} and Corollary \ref{cor24}, we have shown that $\mu_n$ is absolutely continuous with respect to $\mu$ with an $L_{\infty}$ density. Hence, by \cite[Theorem 3.4]{VAD}, we obtain $D_r(\mu_n) \le D_r(\mu).$
%   First, we prove that these results hold when mappings $\{f_i\}{i=1}^N$ are similarity with equal similarity ratio $s_i=s$ for all $i= 1,2, \ldots,N$ and probability vector $\{\rho_i=\rho\}_{i=1}^N$. Observe that, the function $\frac{x}{r+x}$ is strictly increasing and the function $ x \to \sum_{i=1}^N (\rho_i s_i^r)^x$ is strictly decreasing continuous function on $(0, \infty)$. Also, for $x=0$ it is $N$ and it is zero as $x$ approaches $\infty$. Hence, by the intermediate value theorem, there exists a unique $D_r \in (0, \infty)$ such that 
%   $$\sum_{i=1}^N (\rho_i s_i^r)^{\frac{D_r}{r+D_r}}=1.$$
%   The unique number $D_r$ is an upper bound for the quantization dimension, and is equal to the quantization dimension of $\mu$ whenever $\mu$ is a self-similar measure with OSC. Consider $s_i= s$ and $\rho_i= \rho$ then we have,
%   \begin{equation}\label{eq16}
%    N (\rho s^r)^{\frac{D_r(\mu)}{r+D_r(\mu)}}=1.   
%   \end{equation}
% Also, for self-similar measure $\mu_n$ corresponding to the sub-IFS at the $n$-th level, we have 
% \begin{equation}\label{eq17}
%        M (\rho s^r)^{(n+l)\frac{D_r(\mu_n)}{r+D_r(\mu_n)}}=1
%        \implies  M^{-(n+l)} (\rho s^r)^{\frac{D_r(\mu_n)}{r+D_r(\mu_n)}}=1,
% \end{equation}
% where $2 \le M \le N^n$ is a fixed constant. Comparing Equations \ref{eq16} and \ref{eq17}, we obtained $D_r(\mu_n) \le D_r(\mu).$ Also, 
%%%In the above calculation $\rho$ is equal to $\frac{\rho}{\sum_{\eta \in \Upsilon^n}\rho}= \frac{1}{M}$ where $M$ is the number of mappings in The sub-IFS.

  
  
\end{proof}
% In view of the above lemma, we provide two examples below.
% \begin{example}
%     Consider the IFS $\big\{f_1(x) = \frac{x}{2}, f_2(x)= \frac{x}{2}+\frac{1}{2}; \rho_1= \frac{1}{3}, \rho_2= \frac{2}{3}\big\}$ with similarity ratio $s_1=s_2=\frac{1}{2},$ and the sub-IFS at the $3$rd level $\big\{f_{112}, f_{111}, f_{122}; q_{112}= \frac{\rho_{112}}{\rho_{112}+\rho_{111}+\rho_{122}},q_{111}= \frac{\rho_{111}}{\rho_{112}+\rho_{111}+\rho_{122}},q_{122}= \frac{\rho_{122}}{\rho_{112}+\rho_{111}+\rho_{122}}\big\}$ with similarity ratio $s_{111}= s_{122}=s_{122}=\frac{1}{8}.$
%     Clearly, these IFSs satisfy the OSC, then by Graf-Luschgy \cite{GL1}, the quantization dimension $D_2(\mu)$ and $D_2(\mu_n)$ of order $2$ of $\mu$ and $\mu_n$, respectively, are given by the unique solution of the following expressions.
%     \[
% [\rho_1 s_1^2]^{\frac{D_2(\mu)}{2+ D_2(\mu)}} +[\rho_2 s_2^2]^{\frac{D_2(\mu)}{2+ D_2(\mu)}}=1,
% \]
%     \[
%     [q_{111} s_{111}^2]^{\frac{D_2(\mu_3)}{2+ D_2(\mu_3)}}+[q_{112} s_{112}^2]^{\frac{D_2(\mu_3)}{2+ D_2(\mu_3)}}+[q_{122} s_{122}^2]^{\frac{D_2(\mu_3)}{2+ D_2(\mu_3)}}=1
%     \]
%     By putting the respective values and solving these expressions, we get 
%     $D_2(\mu) = 0.9724$ 
%     and $D_2(\mu_3)= 0.513324$, respectively. 
%  Now, 
% \end{example}
% \begin{example}
%     Consider IFS $\big\{f_1(x)= \frac{x}{3}, f_2(x)= \frac{x}{3}+\frac{2}{3}; \rho_1= \frac{1}{7}, \rho_2= \frac{6}{7}\big\},$ with similarity ratio $s_1=s_2=\frac{1}{3}.$ We write $\Upsilon^3= \{111, 112, 121, 122\}$ and the sub-IFS at the $3$rd level 
% $\big\{ f_i; q_i : i \in \Upsilon^3\big\}$
%  with similarity ratio 
% $ s_i =\frac{1}{27}$ for each $i \in \Upsilon^3$,
% where the probability vector is defined by $q_i = \frac{\rho_i}{\sum_{i \in \Upsilon^3} \rho_i}.$ Clearly, these IFSs satisfy the SSC. Then, by \cite[Theorem 7.4.1]{Fraser}
% $$\dim_L(\mu)= \min_{i\in \{1,2\}} \frac{\log \rho_i}{\log s_i} \quad \dim_L(\mu_3)= \min_{i \in \Upsilon^3} \frac{\log q_i}{\log s_i.}$$
% By substituting the respective values, we obtained
% $\dim_L(\mu)\approx 0.14031 $ and $\dim_L(\mu_3) \approx 0.09356.$
% \end{example}

\begin{remark}
We have some IFSs that does not satisfy the separation properties but the sub-IFS at some level satisfy the separation properties. Let $\mathcal{I}=\{\mathbb{R}:f_1,f_2,f_3\}$ be an IFS, where $f_1(x)=\frac{x}{4},~f_2(x)= \frac{x}{4}+t,~f_3(x)=\frac{x}{4}+\frac{3}{4}$ for a fixed constant
$ \frac{1}{16} < t < \frac{4}{16} ~\text{and}~\frac{8}{16} < t < \frac{11}{16}.$

Then $\{f_{11},f_{21},f_{31}\}$ will satisfy the SSC:
\[
f_{11}([0,1])= \Big[0,\frac{1}{16}\Big];~ f_{21}([0,1])=\Big[t,t+\frac{1}{16}\Big];~ f_{31}([0,1])=\Big[\frac{3}{4},\frac{13}{16}\Big].
\]
But the IFS $\mathcal{I}$ does not satisfy the OSC (equivalently, the SOSC) (see Hochman \cite{MHOCHMAN}).
Let $\mathcal{I}=\{\mathbb{R}:f_1,f_2,f_3\}$ be an IFS, where
\[
f_1(x)=\frac{x}{2},~f_2(x)= \frac{x}{4},~f_3(x)=\frac{x}{2}+\frac{1}{2}.
\]
Then $\dim_H(A)=1=\min\{s,1\}$ but for each $\eta \in \Lambda^*$ the IFS $\{f_{\xi \eta}: \xi \in \Lambda\}$ does not satisfy the SSC.
\end{remark}


We note the following fact, which will be used in the proof of the upcoming theorems.
\begin{note}\label{mn1}
The function $\frac{x}{r+x}$ is strictly increasing function on $(0,\infty),$ for any $r>0.$ Also, it is an easy observation that the function $ x \to \sum_{i=1}^N (\rho_i s_i^r)^x$ is strictly decreasing continuous function on $(0,\infty).$
\end{note}
\begin{theorem}\label{mqd1}
Let $\{\mathbb{R}^m; f_{i},\rho_i: i \in \Lambda\}$ be a WIFS consisting of similitude with the similarity ratio $\{s_i\}_{i=1}^N.$ For some $\xi \in \Lambda^*,$ consider the sub-IFS $:= \{f_{ \eta \xi}: \eta \in \Upsilon^n \},$ satisfying the OSC for each $n \in \mathbb{N}$. Then for any $r>0,$ the quantization dimension of order $r$ of the self-similar measure $\mu = \sum_{i \in \Lambda} \rho_{i} \mu_{\rho} \circ f_{i}^{-1}$ associated with the IFS $\{f_{i},\rho_i\}_{i=1}^N$ is $\min\{\kappa_r,m\},$ where $\kappa_r$ is uniquely determined by the expression $\sum_{i \in \Lambda} (\rho_{i} s_{i}^r)^{\frac{\kappa_r}{r + \kappa_r}}= 1.$
\end{theorem}
\begin{proof}
Given that the sub-IFS $\{f_{ \eta \xi}: \eta \in \Upsilon^n\}$ satisfies the OSC, it follows from Graf-Luschgy \cite{GL1} that the quantization dimension of the self-similar measure 
$\mu_n:= \sum_{\eta \in \Upsilon^n} \varsigma_{\eta} \mu_n \circ f_{ \eta \xi}^{-1},$
where $\mu_n$ is the self-similar measure associated with the sub-IFS and the corresponding probability vector 
$\{\varsigma_{\eta}= \frac{\rho_{ \eta \xi}}{ \sum_{\theta \in \Upsilon^n} \rho_{\theta \xi}} : \eta \in \Upsilon^n\}$ 
is determined by the unique solution $\kappa_{n,r}$ of the equation
$\sum_{\eta \in \Upsilon^n} (\varsigma_{\eta} s_{\xi \eta}^r)^{\frac{\kappa_{n,r}}{r + \kappa_{n,r}}}= 1.$
As $D_r(\mu_n)$ may be sufficientely small and tending to zero as $n$ goes to large, we may assume that $\kappa_{n,r}=D_r(\mu_n) \le D_r(\mu_{\rho}).$ For contradiction, we assume that $D_r(\mu_{\rho}) < \kappa_{r}$. Then, using Note \ref{mn1} and $\sum_{\eta \in \Upsilon^n} (\varsigma_{\eta} s_{\xi \eta}^r)^{\frac{\kappa_{n,r}}{r + \kappa_{n,r}}}=1,$ we have 
\begin{align*}
s_{\xi}^{-\frac{r \kappa_{n,r}}{r + \kappa_{n,r}}} 
&= \sum_{\eta \in \Upsilon^n} (\varsigma_{\eta}  s_{ \eta}^r)^{\frac{\kappa_{n,r}}{r + \kappa_{n,r}}}  \\
&\ge \bigg(\sum_{\eta \in \Upsilon^n} \rho_{ \eta}\bigg)^{\frac{-D_r(\mu)}{r+ D_r(\mu)}}\rho_{\xi}^{\frac{D_r(\mu)}{r+ D_r(\mu)}} \sum_{\eta \in \Upsilon^n} (\rho_{ \eta}  s_{ \eta}^r)^{\frac{D_r(\mu)}{r + D_r(\mu)}} \\
& = \bigg(\sum_{\eta \in \Upsilon^n} \rho_{\xi \eta}\bigg)^{\frac{-D_r(\mu)}{r+ D_r(\mu)}}\rho_{\xi}^{\frac{D_r(\mu)}{r+ D_r(\mu)}} \sum_{\eta \in \Upsilon^n} (\rho_{ \eta}  s_{ \eta}^r)^{\frac{D_r(\mu)}{r + D_r(\mu)}}(\rho_{ \eta}   s_{ \eta}^r)^{\frac{-\kappa_{r}}{r + \kappa_{r}}} (\rho_{ \eta}  s_{ \eta}^r)^{\frac{\kappa_{r}}{r + \kappa_{r}}} \\
% &\ge \bigg(\sum_{\eta \in \Upsilon^n} \rho_{\xi \eta}\bigg)^{\frac{-D_r(\mu)}{r+ D_r(\mu)}} \rho_{\xi}^{\frac{D_r(\mu)}{r+ D_r(\mu)}} \sum_{\eta \in \Upsilon^n} (\rho_{ \eta}   s_{ \eta}^r)^{\frac{D_r(\mu)}{r + D_r(\mu)}}(\varsigma_{\eta}  s_{ \eta}^r)^{\frac{-\kappa_{r}}{r + \kappa_{r}}} (\varsigma_{\eta}  s_{ \eta}^r)^{\frac{\kappa_{r}}{r + \kappa_{r}}} \\
&= \bigg(\sum_{\eta \in \Upsilon^n} \rho_{\eta}\bigg)^{\frac{-D_r(\mu)}{r+ D_r(\mu)}}\sum_{\eta \in \Upsilon^n} (\rho_{\eta}  s_{ \eta}^r)^{\frac{D_r(\mu) }{r + D_r(\mu)}- \frac{\kappa_{r}}{r + \kappa_{r}}} (\rho_{\eta}  s_{ \eta}^r)^{\frac{\kappa_{r}}{r + \kappa_{r}}}\\
& \ge  \sum_{\eta \in \Upsilon^n} (\rho_{\eta}  s_{\min}^{nr})^{\frac{D_r(\mu) }{r + D_r(\mu)}- \frac{\kappa_{r}}{r + \kappa_{r}}} (\rho_{\eta}  s_{ \eta}^r)^{\frac{\kappa_{r}}{r + \kappa_{r}}} \\
& \ge  s_{\min}^{nr. \bigg({\frac{D_r(\mu) }{r + D_r(\mu)}- \frac{\kappa_{r}}{r + \kappa_{r}}}\bigg)}  \sum_{\eta \in \Upsilon^n} \rho_{\eta}^{\frac{D_r(\mu) }{r + D_r(\mu)}- \frac{\kappa_{r}}{r + \kappa_{r}}} (\rho_{\eta}  s_{ \eta}^r)^{\frac{\kappa_{r}}{r + \kappa_{r}}} \\ 
 & \ge s_{\min}^{nr. \bigg({\frac{D_r(\mu) }{r + D_r(\mu)}- \frac{\kappa_{r}}{r + \kappa_{r}}}\bigg)}  \sum_{\eta \in \Upsilon^n} \rho_{\eta}^{\frac{D_r(\mu) }{r + D_r(\mu)}} (s_{ \eta}^r)^{\frac{\kappa_{r}}{r + \kappa_{r}}}
\\
&  \ge  s_{\min}^{nr. \bigg({\frac{D_r(\mu) }{r + D_r(\mu)}- \frac{\kappa_{r}}{r + \kappa_{r}}}\bigg)} \sum_{\eta \in \Upsilon^n}( \rho_{\eta} s_{ \eta}^r)^{\frac{\kappa_{r}}{r + \kappa_{r}}}\\ & \ge s_{\min}^{nr. \bigg({\frac{D_r(\mu) }{r + D_r(\mu)}- \frac{\kappa_{r}}{r + \kappa_{r}}}\bigg)}.
 \end{align*}
Since $D_r(\mu)- \kappa_r < 0$ as $n$ approaches infinity, leading to a contradiction, this completes the assertion.
\end{proof}
\begin{remark}
    By \cite[Corollary 11.4]{GL1}, for $r \in [1, \infty),$ the quantization dimension function is a non-decreasing function in $r$. One can observe that this result holds for $r \in [0, \infty)$ due to the fact that for $0 \le p \le q < \infty, ~~ \|x\|_{p} \le \|x\|_{q}.$
\end{remark}
\begin{remark}
    Let $\mu$ be a self similar measure associated with an WIFS $(\{f_i\}_{i=1}^N;\\ \{\rho_i\}_{i=1}^N).$ Then, by \cite[Lemma 3.5]{GL3} $$\lim_{r \to 0^+} D_r(\mu)= D_{0}(\mu).$$
  \end{remark}
  \subsection{Approximation of measures in terms of quantization dimension}
 \begin{lemma}\label{lem10}
     Let $\mu$ and $\nu$ be two finite measures. Then 
     $$D_{0}(\mu +\nu) = \max\{D_{0}(\mu) , D_{0}(\nu)\}.$$
 \end{lemma}
 \begin{proof}
     By the definition of the $n$th quantization error of a measure with respect to the geometric mean error, we have 
     \begin{equation}
      \begin{aligned}
          \log e_{n}(\mu+\nu) &= \inf \bigg\{ \int \log d(x,A) d(\mu+\nu)(x) : A \subset \mathbb{R}^d, \text{Card}(A) \le n\bigg\} \\& \ge 
\inf \bigg\{ \int \log d(x,A) d(\mu)(x)\bigg\} + \inf \bigg\{ \int \log d(x,A) d(\nu)(x)\bigg\} \\& = \log e_n(\mu) + \log e_{0}(\nu).
  \end{aligned}   
     \end{equation}
     This depicts that $D_{0}(\mu + \nu) \ge \max\{D_{0}(\mu), D_{0}(\nu)\}.$ Now, let $t > \max\{D_{0}(\mu), D_{0}(\nu)\},$ then by \cite[Proposition 4.3]{GL3}, we have 
     $$\lim_{n \to \infty}(\log n + t \log e_{n}(\mu)) = +\infty \quad \lim_{n \to \infty}(\log n + t \log e_{n}(\nu)) = +\infty.$$
     Let $k \in \mathbb{R}$ be any arbitrary number, then there exist natural numbers $N_1(\mu,k)$ and $N_2(\nu,k)$ such that 
     $$\log n  + t \log e_{n}(\mu) > k \quad \text{and} \quad \log n  + t \log e_{n}(\mu) > k \quad \forall n \ge N = \max\{N_1,N_2\}.$$
     Also, by the definition of the quantization error of a measure with respect to the geometric mean error, there exist sets $A(\mu,n,k)$ and $B(\nu,n,k)$ such that 
     \begin{equation*}
         \begin{aligned}
            &\log n  + t  \int \log d(x,A(\mu,n,k)) d(\mu)(x) > k, \quad \text{and} \\& \log n  + t \int \log d(x,B(\nu,n,k)) d(\nu)(x) > k \quad \forall n \ge N. 
         \end{aligned}
     \end{equation*}
  Since Card$(A(\mu,n,k) \le n)$ and Card$(B(\nu,n,k) \le n),$ it follows that Card$(A \cup B) \le n^2,$ and we have $$2 \log n + t \int \log d(x, A\cup B) d (\mu)(x) > k.$$   
  Now, for every $n > N,$ we have 
  \begin{equation*}
      \begin{aligned}
          2 \log n + t \int \log d(x, A \cup B) d(\mu + \nu)(x) & =  \log n + t \int \log d(x, A \cup B) d(\mu)(x) + \\& \log n + t \int \log d(x, A \cup B) d(\nu)(x) > k
      \end{aligned}
  \end{equation*}
     Since $k \in \mathbb{R}$ is arbitrary, we have $\lim_{n \to \infty}(\log n+ t \log e_{n}(\mu + \nu )) = +\infty$. It follows that $D_{0}(\mu + \nu) < t.$ Further, since $t > \max\{D_{0}(\mu), D_{0}(\nu)\}$ is arbitrary, it follows that $D_{0}(\mu + \nu) \le \max\{D_{0}(\mu), D_{0}(\nu)\}.$ This completes the assertion.
 \end{proof}
\begin{lemma}\label{lem09}
 Let $\mu$ be a finite Borel probability measure. Then for a fixed $\nu \in \mathcal{L}(\mathbb{R}^m),$ 
 $$D_{0}(\mu * \nu) = D_{0}(\mu).$$
\end{lemma}
\begin{proof}
    Following a similar approach to \cite[Lemma 3.10]{VAD}, we obtain the required result.
\end{proof}
Using the above lemmas, we will approximate the set $\Omega(\mathbb{R}^m)$ by several of its subclasses with respect to the quantization dimension.

% \begin{proof}
%     % As $r\to 0^+$, $D_r(\mu) \to \dim_H(\mu)$ for any self-similar measure (OSC). 
%     Since $D_r$ is non-decreasing function in $r,$ we have $D_0(\mu) \le D_r(\mu)$ for all $r>0.$ Thus, we need to show that $\limsup_{r \to 0^+} D_r(\mu) \le D_0(\mu).$
% \end{proof}

\begin{theorem}
    The sets $\Omega_0({\mathbb{R}^m})=\{\mu \in \Omega(\mathbb{R}^m): \lim_{r \to 0^+}D_r(\mu) =D_0(\mu)= \dim_H(\mu)\}$ and $\Omega_{0>}({\mathbb{R}^m})=\{\mu \in \Omega(\mathbb{R}^m): \lim_{r \to 0^+}D_r(\mu) =D_0(\mu)>\dim_H(\mu)\}$  are dense in $\Omega(\mathbb{R}^m).$
\end{theorem}
\begin{proof}
Since By \cite[Remark 4.6]{GL3}, the set $\Omega_0({\mathbb{R}^m})$ is non-empty. Let $\mu \in \Omega(\mathbb{R}^m)$ be any Borel probability measure and $\epsilon >0,$ then by the fact that $\mathcal{L}(\mathbb{R}^m)$ is dense in $\Omega(\mathbb{R}^m),$ there exists a measure $\nu \in \mathcal{L}(\mathbb{R}^m)$ such that
$$d_L(\mu, \nu) < \frac{\epsilon}{2}.$$ Further, by Lemma \ref{lem09} and the similar technique followed in \cite[Theorem 3.16]{VAD}, we have there exists a compactly supported measure such that $d_L(\nu, \lambda) < \frac{\epsilon}{2},$ where $\lambda = \nu * \lambda_1$ for some $\lambda_1 \in \Omega_0({\mathbb{R}^m}).$ in this way, we have $$d_L(\mu, \lambda_1) \le d_L(\mu,\nu)+ d_L(\nu,\lambda) < \epsilon.$$
This proves that the set $\Omega_0({\mathbb{R}^m})$ is a dense subset of $\Omega(\mathbb{R}^m).$
Again, by \cite[Theorem 4.3]{GL3} the set $\Omega_{0>}({\mathbb{R}^m})$ is non-empty. Since the proof is similar, we omit it.

\end{proof}















% In order to prove the upcoming theorem, we recall the code-space structure on the IFS and the associated sub-IFS at the $n$th level.\\
% Let $(\{f_i\}_{i=1}^N\{\rho_i\}_{i=1}^N)$ be a WIFS. We denote $\Sigma^n:= \{1,2,\ldots,N\}^n=\{\eta_1,\eta_2,\ldots,\eta_n: 1 \le \eta_i \le N\}$, $\Sigma^* =\bigcup_{k \in \mathbb{N}} \Sigma^k$, and $\Sigma^*= \{1,2,\ldots,N\}^{\mathbb{N}}.$ The code space $\Sigma$ endowed with the metric $d(\eta, \omega) := \alpha^k,$ where $k = \min_{i}\{i : \eta_i \neq \omega_i\}$ and $\alpha \in (0,1).$
% The 
% The Bernoulli measure $\mu$ on 
% \begin{theorem}\label{mld1}
% Let $\{\mathbb{R}^m;\Phi_{i},\rho_i: i \in \Lambda\}$ be a WIFS consisting of a similarity with the similarity ratio $\{s_i\}_{i=1}^N.$ For some $\xi \in \Lambda^*,$ consider the IFS $:= \{\Phi_{ \eta \xi}: \eta \in \Upsilon^n \},$ satisfying the SSC for each $n \in \mathbb{N}$. Then, the lower dimension of the self-similar measure $\mu_{\rho} = \sum_{i \in \Lambda} \rho_{i} \mu_{\rho} \circ \Phi_{i}^{-1}$ associated with the IFS $\{\Phi_{i},\rho_i\}_{i=1}^N$ is.....
% $$\dim_L(\mu_\rho)= \min_{i \in \Lambda} \frac{\log (\rho_i)}{\log (s_i)}$$
% \end{theorem}
% \begin{proof}
% The sub-IFS $\{\Phi_{ \eta \xi}: \eta \in \Upsilon^n\}$ satisfies the SSC, from \cite[Theorem 7.4.1]{Fraser}, we deduced that the lower dimension of the self-similar measure $\mu_n = \sum_{\eta \in \Upsilon^n} \varsigma_{\eta} \mu_{n} \circ \Phi_{\eta \xi}^{-1}$, where $\mu_n$ is the self-similar measure associated with the sub-IFS and the corresponding probability vector $\{ \varsigma_\eta = \frac{\rho_{\eta \xi}}{ \sum_{\theta \in \Upsilon^n} \rho_{\theta \xi}} : \eta \in \Upsilon^n\}$ is determined by $\dim_L(\mu_n) =\min_{\eta \in \Upsilon^n} \frac{\log (\varsigma_\eta)}{\log (s_{\eta\xi})}$.

% Clearly, $\mu_n < \mu$, where $\mu$ is a self-similar measure associated with $\Lambda$. Therefore, $ \dim_S(\mu_n) = \dim_{L}(\mu_n) \leq \dim_{L}(\mu_{\rho})$.

% we have 
% \begin{align*}
% \dim_{L}(\mu) &\ge \dim_S(\mu_n) = \min_{\eta \in \Upsilon^n} \frac{\log \varsigma_\eta }{ \log s_{\eta\xi}} \\
%  &= \min_{\eta \in \Upsilon^n} \frac{\log \varsigma_\eta }{ \log (s_{\eta}s_{\xi})}\\
%   &= \min_{\eta \in \Upsilon^n} \frac{\log \frac{\rho_{\eta \xi}}{\sum_{\theta \in \Upsilon^n} \rho_{\theta \xi}}}{ \log s_{\eta} + \log s_{\xi}} \\
  % &= \min_{\eta \in \Upsilon^n} \frac{\log \rho_{\eta \xi} - \log {\sum_{\theta \in \Upsilon^n} \rho_{\theta \xi}}}{ \log s_{\eta} + \log s_{\xi}}\\
  % &= \min_{\eta \in \Upsilon^n} \frac{\log \rho_{\eta} \rho_{\xi} - \log {\sum_{\theta \in \Upsilon^n} \rho_{\theta \xi}}}{ \log s_{\eta} + \log s_{\xi}}\\
  % &= \min_{\eta \in \Upsilon^n} \frac{\log \rho_{\eta} +\log\rho_{\xi} - \log {\sum_{\theta \in \Upsilon^n} \rho_{\theta \xi}}}{ \log s_{\eta} + \log s_{\xi}}
% \end{align*}
% \end{proof}




% Recall that a set is said to be $G_\delta$ if it can be written as a countable intersection of open sets.
% % \begin{theorem}
% %     The set $\{\mu \in \Omega(\mathbb{R}^m): D_r(\mu)=0\}$ is $G_\delta$-dense subset of $\Omega(\mathbb{R}^m)$.
% % \end{theorem}
% % \begin{proof}
% %     let...
% % \end{proof}
% Next, we will focus on the convolutional decomposition of measures in terms of dimension-preserving.

% \begin{theorem}
% For $\mu \in \Omega(\mathbb{R}^m)$ and $\beta \in [0,m]$, there exist two measures $\nu$ and $\lambda$ such that $\mu=\nu * \lambda$ and $D_r(\nu)=D_r(\lambda) =\beta.$ 
% \end{theorem}
% \begin{proof}
%     let...
% \end{proof}


% \begin{theorem}\label{densethm}
% The set  $\mathcal{S}_{0}:=\{ \mu \in  \Omega(\mathbb{R}^m): D_r(\mu)  = 0 \}$ is a $G_{\delta}$-dense (generic) subset of $\Omega(\mathbb{R}^m).$
% \end{theorem}
% \begin{proof}
% Let $\beta >0$. Define $\mathcal{B}_{\beta}=\{\mu \in  \Omega(\mathbb{R}^m): \lim_{n \to \infty}n e_{n,r}(\mu)^{\beta}=0\},$ where $e_{n,r}(\mu)= V_{n,r}^{\frac{1}{r}}(\mu).$ Using the fact that $ne_{n,r}(\mu)^{\beta}$ is a continuous function, observe that $\mathcal{B}_{\beta}$ is $G_{\delta}$-set. Now, $\mathcal{S}_{0}= \cap_{\beta>0} \mathcal{B}_{\beta}.$ Using the fact that countable intersection of $G_{\delta}$-sets is a $G_{\delta}$-set, $\mathcal{S}_{0}$ is a $G_{\delta}$-set. This, with the previous theorem, establishes the result.
% \end{proof}
 % \begin{theorem}\label{densethm1}
 % The set $\mathcal{S}_{m}:=\{ \mu \in  \Omega(\mathbb{R}^m): D_r(\mu)  = m \}$ is a prevalent subset of $\Omega(\mathbb{R}^m).$
 % \end{theorem}
 % \begin{proof}
 % Let...(Not proved)...
 % \end{proof}
We introduce the total variation metric (\cite{Parth}) on $\Omega( \mathbb{R}^m)$ as follows:
\[
\|\mu - \nu\|_{TV}=\sup\{|\mu(A)- \nu (A)|: A ~\text{is a Borel set}\}.
\]

\begin{lemma}\label{lemcon1}
            
Let $\lambda \in \Omega( \mathbb{R}^m)$ be a fixed measure and $\big(\mu_n\big)$ be a sequence of absolutely continuous (with respect to $\lambda$) measures in $\Omega( \mathbb{R}^m).$ If $(\frac{d\mu_n}{d\lambda})$ converges in $\mathcal{L}^1(\lambda)$-norm then $\big(\mu_n\big)$ converges to a measure $\mu$, and
$$
\frac{d\mu}{d\lambda}=\lim_{n \to \infty }\frac{d\mu_n}{d\lambda}.
$$
\end{lemma}
\begin{proof}
Let $g= \lim_{n \to \infty} \frac{d\mu_n}{d\lambda}.$ Then a measure $\mu$ defined by $$\mu(A)= \int_A g d\lambda, ~~ ~~~~ A \subset \mathbb{R}^n,$$
is in $ \Omega( \mathbb{R}^m).$ Further, we obtain 
\begin{equation*}
\begin{aligned}
d_L(\mu_n,\mu)& =\sup_{\te{Lip}(f)\leq 1} \Big\{\Big|\int_{\mathbb{R}^m} f d\mu_n -\int_{\mathbb{R}^m} fd\mu \Big|\Big\} \\ & = \sup_{\te{Lip}(f)\leq 1} \Big\{\Big|\int_{\mathbb{R}^m} f  \frac{d\mu_n}{d\lambda} d\lambda -\int_{\mathbb{R}^m} fgd\lambda \Big|\Big\} \\ & \le  \sup_{\te{Lip}(f)\leq 1} \Big\{\int_{\mathbb{R}^m} \Big|f  \frac{d\mu_n}{d\lambda} -fg \Big|d\lambda\Big\} \\ & \le  \sup_{\te{Lip}(f)\leq 1} \|f\|_{\infty} \Big\{\int_{\mathbb{R}^m} \Big|  \frac{d\mu_n}{d\lambda} -g \Big|d\lambda\Big\} \\ & \le   \int_{\mathbb{R}^m} \Big|  \frac{d\mu_n}{d\lambda} -g \Big|d\lambda.
\end{aligned}
\end{equation*}
Since $\frac{d\mu_n}{d\lambda} \to g$ in $\mathcal{L}^1(\lambda)$-norm, we have $d_L(\mu_n,\mu) \to 0$ as $n \to \infty,$ hence the proof.
\end{proof}
\begin{remark}
Under the hypothesis of the above theorem, we also have 
\begin{equation*}
\begin{aligned}
\|\mu_n-\mu\|_{TV}& =\sup\{|\mu_n(A)- \mu (A)|: A ~\text{is a Borel set}\}\\ & = \sup\Big\{\Big|\int_{A}   \frac{d\mu_n}{d\lambda} d\lambda -\int_{A} gd\lambda \Big| : A ~\text{is a Borel set}\Big\} \\ & \le  \sup \Big\{\int_{A} \Big|  \frac{d\mu_n}{d\lambda} -g \Big|d\lambda:  A ~\text{is a Borel set}\Big\} \\ & \le   \int_{\mathbb{R}^m} \Big|  \frac{d\mu_n}{d\lambda} -g \Big|d\lambda.
\end{aligned}
\end{equation*}
From the above, we infer that $\|\mu_n-\mu\|_{TV} \to 0$ as $n \to \infty.$
\end{remark}
\begin{theorem}\label{approxthm1}
Let  $\lambda \in \Omega( \mathbb{R}^m).$
Suppose $\mu$ is a absolutely continuous measure with respect to $\lambda$ such that $\dim_H \Big(Gr\big(\frac{d\mu}{d\lambda}\big)\Big) =\alpha$ for some $m \le \alpha \le m+1.$ Then there exists a sequence of absolutely continuous measures $(\mu_n) \in \Omega( \mathbb{R}^m)$ with the following conditions: $$\dim_H\Big(Gr\big(\frac{d\mu_n}{d\lambda}\big)\Big) =\alpha ~~\text{and}~~ \mu =\lim_{n \to \infty}\mu_n.$$  
\end{theorem}
\begin{proof}
In the light of an analogous result to Theorem \ref{lemcon1}, we may have a sequence of integrable (with respect to $\lambda$) functions $(g_n)$ such that $\dim_H\big(Gr(g_n)\big) =\alpha$ and  $g_n \to \frac{d\mu}{d\lambda}$ in $\mathcal{L}^1(\lambda)$-norm. 
Now, one defines a sequence of measures $(\mu_n)$ as follows $$\mu_n(A)=\int_A g_nd\lambda, ~~~~~~ A \subset \mathbb{R}^m.$$ Then, $\frac{d\mu_n}{d\lambda}=g_n$ and $\Big(\frac{d\mu_n}{d\lambda}\Big)$ converges to $\frac{d\mu}{d\lambda}.$ Lemma \ref{lemcon1} yields that the sequence $(\mu_n)$ converges to $\mu$ such that $\dim_H\Big(Gr\big(\frac{d\mu_n}{d\lambda}\big)\Big) =\alpha.$ Hence the assertion is proved. 
\end{proof}
% \section{Decomposition of measures in terms of quantization dimension}
% The next theorem is motivated by Theorem $1.2$ of \cite{Liu}.
% \begin{theorem}\label{maindeco}
% Let $0\le \alpha \le m$ and $\mu \in \Omega( \mathbb{R}^m).$ Then,  there exist Borel probability measures $\nu,\lambda$ such that $\mu=\nu * \lambda$ and $D_r(\nu)=D_r(\lambda)=\alpha.$
% \end{theorem}
% \begin{proof}
% I think this theorem holds, in general:

% \end{proof}


% The next theorem is motivated by Theorem $1.2$ of \cite{Liu}.
% \begin{theorem}\label{maindeco}
% Let $0\le \alpha \le m$, and $\mu$ be a finite Borel measure such that $D_r(\mu) \le \alpha.$ Then, there exist finite Borel measures $\nu,\lambda$ such that $\mu=\nu+\lambda$ and $D_r(\nu)=D_r(\lambda)=\alpha.$
% \end{theorem}
% \begin{proof}
% Let $D_r(\mu)=\alpha.$ Choose $\nu=\lambda= \frac{1}{2}\mu.$ By Lemma \ref{lipdim}, we have $D_r(\nu)=D_r(\lambda)=\alpha,$ hence the result whenever $D_r(\mu)=\alpha.$ 
% \end{proof}
% \section{Approximation of measures in terms of generalized Hausdorff dimension}
%  In this section, we will approximate the space $\Omega(\mathbb{R})^n$ with respect to the generalized $\psi_t$-Hausdorff dimension.
%  .......




%  \begin{theorem}
%   For any $0 \le \beta \le 1,$ the set $\mathcal{D}_1:= \{\mu \in \Omega(\mathbb{R}^n) : \dim_{H}^{\Psi_t}(\mu)= \beta~\text{and}~\mu << \mathcal{L}^1 \}$ is dense in $\Omega(\mathbb{R})$.
%   \end{theorem}
% \begin{proof}
%   ......  
% \end{proof}
% \begin{theorem}
%    For any $0 \le \beta \le 1,$ the set $\mathcal{D}_2= \{\mu \in \Omega(\mathbb{R}^n) : \dim_{H}^{\Psi_t}(\mu)= \beta~\text{and}~\mu \bot \mathcal{L}^1\}$ is dense in $\Omega(\mathbb{R}).$ 
% \end{theorem}
% \begin{proof}
%   ......  
% \end{proof}


% \begin{theorem}
%     The set $\mathcal{S}_0 := \{ A \in \mathcal{K}(\mathbb{R}) : \mathcal{H}^{\Psi_t}(A)=0$ is dense in $\mathcal{K}(\mathbb{R}).$
% \end{theorem}
% \begin{proof}
%  ......   
% \end{proof}
% \begin{theorem}
%     The set $\mathcal{S}_{\infty}:=\{A \in \mathcal{K}(\mathbb{R}): \mathcal{H}^{\Psi_t}(A)= \infty\}$ is dense in $\mathcal{K}(\mathbb{R}).$
% \end{theorem}
% \begin{proof}
%     .....
% \end{proof}
% \begin{theorem}
%  The set $\mathcal{S}_f:= \{A \in \mathcal{K}(\mathbb{R}): 0<\mathcal{H}^{\Psi_t}(A) < + \infty\}$ is dense in $\mathcal{K}(\mathbb{R}).$  
% \end{theorem}
% \begin{proof}
%     ......
% \end{proof}
\section*{Acknowledgements}
 The first author is supported by the SEED grant project of IIIT Allahabad. The second author is financially supported by the Ministry of Education at the IIIT Allahabad. Some parts of the paper have been presented in lectures delivered by the first author at the workshop ``Course on Exponential separation and $L^q$ dimension" organized by IIIT Allahabad on April 23, 2025.


\bibliographystyle{amsalpha}
\begin{thebibliography}{A}

% \bibitem{AV1} E. Agrawal, S. Verma, Dimension preserving approximation and estimation: Fractal surfaces and Riemann-Liouville fractional integrals, Contemporary Mathematics (AMS), CONM16505, (2025) to appear.
% \bibitem{AV2} E. Agrawal, S. Verma, Dimension preserving set-valued approximation and decomposition via metric sum, (2025)
% arXiv preprint arXiv:2504.11356.
% \bibitem{AV3} E. Agrawal, S. Verma, Dimensions and stability of invariant measures supported on fractal surfaces, Discrete and Continuous Dynamical Systems-S, (2024) 1-27, doi: 10.3934/dcdss.2024159.
% \bibitem{Arc} N. Arcozzi, A. Monguzzi, M. Salvatori, Ahlfors regular spaces have regular subspaces of any dimension, Riv. Math. Univ. Parma. 14(1) (2023) 19-32.

%  \bibitem{Aubin} J. P. Aubin H. Frankowska, Set-Valued Analysis, Birkh\"{a}user, Boston, 1990.

% \bibitem{BG1} C. Bandt, S. Graf, Self-similar sets VII. A characterization of self-similar fractals with positive Hausdorff measure, Proc. Amer. Math. Soc. 114(4) (1992), 995-1001.

\bibitem{BV1} B. B\'ar\'any, M. Verma, Hausdorff dimension of self-similar measures and sets with common fixed point structure, arXiv preprint, (2025) arXiv:2507.05835.
\bibitem{B} M. F. Barnsley, Fractals Everywhere, 2nd edition, Academic Press, Boston, 1993.


% \bibitem{Bayart} F. Bayart, Y. Heurteaux,
% On the Hausdorff dimension of graphs of prevalent continuous functions on compact sets,
% Further Developments in Fractals and Related Fields, Birkhäuser, Boston (2013) 25-34.


% \bibitem{Bed} T. Bedford, Crinkly curves, Markov partitions and box dimensions in self-similar sets, PhD thesis, University of Warwick, 1984.


\bibitem{Bill} P. Billingsley, Convergence of Probability Measures, Wiley \& Sons, New York- London, 1968.
% \bibitem{SS3} S. Chandra, S. Abbas, On fractal dimensions of fractal functions using function spaces, Bull. Aust. Math. Soc. 106 (2022) 470-480.
% \bibitem{Chen} C. Chen, M. Wu, W. Wu, Accessible values for the Assouad and lower dimensions of subsets, Real Anal. Exchange 45(1) (2020) 85-100. 
% \bibitem{CM} R. Cawley, R. D. Mauldin, Multifractal decompositions of Moran fractals, Adv. Math. 92 (2)
% (1992) 196–236.
% \bibitem{Conway} J. B. Conway, \textit{A course in functional analysis}, 96 Springer, 2019.

% \bibitem{DGSH} A. Deliu, J. S. Geronimo, R. Shonkwiler, D. Hardin, \textit{Dimensions associated with recurrent self-similar sets}, Proc. Cambridge Philos. Soc. \textbf{110}(2) (1991), 327--336. 

% \bibitem{DV1} S. Dubey, S. Verma, Fractal dimension for a class of Inhomogeneous graph-directed attractors, Discrete and Continuous Dynamical Systems-S, (2024) doi: 10.3934/dcdss.2024158.
% \bibitem{Fal31} K. J. Falconer,  Techniques in Fractal Geometry, John Wiley and Sons, 1997.

\bibitem{Fal} K. J. Falconer, Fractal Geometry: Mathematical Foundations and Applications, John Wiley Sons Inc., New York, 1999.

\bibitem{Feng} D. J. Feng, N. T. Nguyen, T. Wang, Convolutions of equicontractive self-similar measures on the line, Illinois J. Math. 46(4) (2002) 1339-1351.

 \bibitem{Fraser12} J. M. Fraser, Assouad type dimensions and homogeneity of fractals, Trans. Amer. Math. Soc. 366(12) (2014) 6687-6733.

 \bibitem{Fraser} J. M. Fraser, Assouad dimension and fractal geometry, (Vol. 222), Cambridge University Press, 2020.

% \bibitem{Fraser1} J. M. Fraser, Fractal Geometry of Bedford-McMullen Carpets, Lecture Notes in Mathematics, Vol 2290. Springer, 2021. 
    % \bibitem {Fal} K. J. Falconer, Fractal Geometry: Mathematical Foundations and Applications, John Wiley Sons Inc., New York, 1999.
    
    \bibitem{GL1} S. Graf, H. Luschgy, Foundations of quantization for probability distributions, Lecture Notes in
    Mathematics 1730, Springer, Berlin, 2000.

    
    %\bibitem{GL2}  S. Graf, H. Luschgy, The Quantization dimension of self-similar probabilities, Math. Nachr. 241
    % (2002) 103-109.

    \bibitem{GL3} S. Graf, H. Luschgy, Quantization for probability measures with respect to the geometric mean error, In Mathematical Proceedings of the Cambridge Philosophical Society 136(3) (2004) 687-717.
    % \bibitem{Gurubachan} Gurubachan, V. V. M. S. Chandramouli, S. Verma, Fractal dimension of $\alpha$-fractal functions without endpoint conditions, Mediterr. J. Math. \textbf{21} (2024), 71.
% \bibitem{Hare} K. Hare, F. Mendivil, L. Zuberman, Measures with specified support and arbitrary Assouad dimensions, Proceedings of the American Mathematical Society 148(9) (2020) 3881-3895.
 \bibitem{MHOCHMAN} M. Hochman, On self-similar sets with overlaps and inverse theorems for entropy, Ann. Math. 180(2) (2014) 773-822.

\bibitem{YH} Y. Heurteaux, Dimension of measures: the probabilistic approach, Publ. Mat. 51(2) (2007) 243-290.
 
\bibitem{H} J. E. Hutchinson, Fractals and self similarity, Indiana Uni. Math. J. 30(5) (1981) 713-747.
 
  % \bibitem{AL} A. K\"aenm\"aki, J. Lehrb\"ack, Measures with predetermined regularity and inhomogeneous self-similar sets, Ark. Mat. 55(1) (2017) 165-184.

  % \bibitem{AJM} A. Käenmäki, J. Lehrbäck, M. Vuorinen, Dimensions, Whitney covers, and tubular neighborhoods, Indiana Univ. Math. J. 62(6) (2013) 1861-1889. 
  
% \bibitem{KN1} M. Kesseböhmer, A. Niemann,T. Samuel, H. Weyer, Generalised  Kreĭn—Feller operators and gap diffusions via transformations of measure spaces.
% To appear in: ``From Classical Analysis to Analysis on Fractals. A Tribute to Robert Strichartz, Volume 2", Applied and Numerical Harmonic Analysis, Birkhäuser, 2024.


\bibitem{KNZ1} M. Kesseböhmer, A. Niemannn, Quantization dimensions of negative order, (2025), 	arXiv:2405.13387.
\bibitem{KNZ1} M. Kesseböhmer, A. Niemannn, S. Zhu, 
Quantization dimensions of compactly supported probability measures via R\'enyi dimensions,
Trans. Amer. Math. Soc. 376 (2023) 4661-4678.
 \bibitem{KZ1} M. Kesseböhmer, S. Zhu,
Some recent developments in quantization of fractal measures,
 Fractal geometry and stochastics V, Progr. Probab. 70: (2015) 105-120 .
\bibitem{KZ2} M. Kesseböhmer, S. Zhu,
On the quantization for self-affine measures on Bedford-McMullen carpets,
Mathematische Zeitschrift  283(1) (2015) 39-58.
\bibitem{Larman} D. G. Larman, A new theory of dimension, Proc. Lond. Math. Soc. 3(1) (1967) 178-192.
\bibitem{Liu} J. Liu, J. Wu,
     A remark on decomposition of continuous functions, J. Math. Anal. Appl. 401 (2013) 404-406.
 % \bibitem{Mauldin} R. D. Mauldin, S. C. Williams, On the Hausdorff dimension of some graphs, Trans. Amer. Math. Soc. 298(2) (1986) 793-803.
      \bibitem {Mat} P. Mattila, Geometry of Sets and Measures in Euclidean Spaces, Cambridge University Press 1995.
    % \bibitem {Mat2} P. Mattila, M. Mor\'an, J.-M. Rey, Dimension of a measure, Studia Math. 142 (2000) 219-233.
     \bibitem {Mat3} P. Mattila, Fourier analysis and Hausdorff dimension, Cambridge University Press, Cambridge, 2015.
     % \bibitem{MP1} M. Megala, S. A. Prasad, Spectrum of a Self-Affine Measure with Four-Element Digit Set, Fractals 30 (04), (2022) 2250087.
% \bibitem{MP2} M. Megala, S. A. Prasad, Spectrum of self-affine measures on the Sierpinski family, Monatshefte f\"ur Mathematik 204(1) (2024) 157-169.
% \bibitem{MP3} M. Megala, S. A. Prasad, Spectrality of certain self-affine measures, Contemporary Mathematics (AMS), Vol. 825, (2025) 171-180.
      \bibitem {Parth} K. R. Parthasarathy, Probability measures on metric spaces, Academic
      Press, New York -London, 1967.
    
      \bibitem{KPot} K. P\"otzelberger, The quantization dimension of distributions, Math. Proc. Cambridge Philos. Soc. 131 (2001) 507-519.
      







\bibitem{PVV1} A. Priyadarshi, M. Verma, S. Verma, Fractal dimensions of fractal transformations and quantization dimensions for bi-Lipschitz mappings, J. Fractal Geom. 12(1/2) (2025) 1-33.



\bibitem{Schief} A. Schief, Separation properties for self-similar sets, Proc. Amer. Math. Soc. 122(1) (1994) 111-115.

\bibitem{Sh1} P. Shmerkin, On Furstenberg's intersection conjecture, self-similar measures, and the $L^q$ norms of convolutions, Ann. of Math. 189(2) (2019) 319-391.


       % \bibitem{Shaw} A. Shaw, Prevalence, M. Math Dissertation, University of St. Andrews, 2010.
% \bibitem{Sh1} P. Shmerkin, On Furstenberg's intersection conjecture, self-similar measures, and the $L^q$ norms of convolutions, Ann. of Math. 189(2) (2019) 319-391.

% \bibitem{Simon} K. Simon, Overlapping cylinders: the size of a dynamically defined Cantor-set, In Ergodic theory of $Z^d$ actions (Warwick, 1993-1994), volume 228 of London Math. Soc. Lecture Note Ser., pages 259--272. Cambridge University Press, Cambridge, 1996.
% %  \bibitem{BS} B. Solomyak, Fourier decay for self-similar measures, Proc. Amer. Math. Soc. 149(8) (2021) 3277-3291.
% \bibitem{BS1} B. Solomyak, Fourier decay for homogeneous self-affine measures, JJ. Fractal Geom. 9(1) (2022) 193-206.

% \bibitem{Suom} V. Suomala, Intermediate value property for the Assouad dimension of measures, Analysis and Geometry in Metric Spaces 8(1) (2020) 106-113.


\bibitem{PPV1} P. P. Varj\'u, On the dimension of Bernoulli convolutions for all transcendental parameters, Ann. of Math. (2) 189(3) (2019) 1001-1011.

% \bibitem{MV1} M. Verma, Dimensional approximation of inhomogeneous attractor without any separation condition, Bull. Aust. Math. Soc. (2025) 1-11.

% \bibitem{MVP1} Manuj Verma, A. Priyadarshi, A note on the dimensions of difference and distance sets for graphs of functions,  Chaos Solitons and Fractals, 184  (2024), pp. 114986.

\bibitem{VAD} S. Verma, E. Agrawal, S. Dubey, Some results on Lower Assouad and quantization dimensions, arxiv preprint (2025) 1-16, DOI: 10.48550/arXiv.2508.15282.
     % \bibitem{VAM} S. Verma, E. Agrawal, M. Megala, Some remarks on the exponential separation and dimension preserving approximation for sets and measures, (2025), 1-19, arXiv preprint, doi.org/10.48550/arXiv.2508.07692.
% \bibitem{VP1} S. Verma, A. Priyadarshi, Further analysis of Hausdorff dimension and separation conditions, Contemporary Mathematics (AMS), Vol. 825, (2025) 225-239, doi.org/10.1090/conm/825/16517.
 % \bibitem{VM1} S. Verma, P. R. Massopust, Dimension preserving approximation, Aequat. Math. 96(6) (2022) 1233-1247.
 \bibitem{XYZL} X. Yang, Z. Liu, Some remarks on convergence of measures, IOSR Journal of Engineering 2(5) (2012) 1069-1070.
    \bibitem{Zador} P. L. Zador, Asymptotic quantization error of continuous signals and the
     quantization dimension, IEEE Trans. Inform. Theory 28 (1982) 139-149.
    % \bibitem{Zhu} S. Zhu, The lower quantization coefficient of the $F$-conformal measure is positive, Nonlin. Anal. 65(2) (2008) 448-455.
    
    \bibitem{Zhu1} S. Zhu, The quantization for self-conformal measures with respect to the geometric mean error, Nonlinearity 23(11) (2010) 2849.
\bibitem{Zhu2} S. Zhu, Quantization for probability measures, PhD diss., Universität Bremen, 2005.
\bibitem{Zhu3} S. Zhu, Some remarks on absolute continuity and quantization of probability measures, Journal of mathematical analysis and applications 330(2) (2007) 1465-1477.
\end{thebibliography}

\end{document}

%-----------------------------------------------------------------------
% End of article.tex
%-----------------------------------------------------------------------
